\section{Mises en candidature}
\label{mises-en-candidature}

\begin{enumerate}
 \item Les mises en candidature aux postes de l'équipe de direction se feront par vote lors d'une assemblée générale précédant l'assemblée générale annuelle chaque année. Les membres voteront sur une liste de candidatures à présenter pour élection lors de l'assemblée générale annuelle.
 \item Les personnes candidates à tout poste élu devront d'abord être nominées par dix (10) membres des associations membres.
  \begin{enumerate}
   \item Les formulaires de mise en candidature seront distribués aux personnes candidates potentielles par la DGE.
  \end{enumerate}
 \item Aucune personne ne sera candidate à plusieurs postes lors de la même élection.
 \item Les personnes candidates admissibles doivent :
  \begin{enumerate}
   \item être âgées d'au moins 18 ans ;
   \item ne pas avoir été déclarées inaptes en vertu des lois d'une province ou d'un territoire canadien, ou par un tribunal dans une juridiction hors du Canada ;
   \item être une personne physique (une société ne peut être une personne administratrice ou un membre de l'équipe de direction) ;
   \item ne pas avoir le statut de failli ; et
   \item consentir par écrit à siéger en tant que personne administratrice et/ou membre de l'équipe de direction, si élue.
  \end{enumerate}
 \item Le vote pour déterminer la liste de candidatures se déroulera de la manière décrite à la section 8 de ces statuts électoraux.
 \item La liste de candidatures sélectionnée sera soumise à l'élection lors de l'assemblée générale annuelle.
 \item Si une personne candidate nominée se retire, devient inadmissible ou est incapable de se présenter à l'élection avant l'assemblée générale annuelle :
  \begin{enumerate}
   \item le poste de cette personne candidate sur la liste de candidatures deviendra vacant ;
   \item le CA ne substituera pas de sa propre autorité une autre personne candidate ;
   \item le CA peut convoquer une autre assemblée de mise en candidature si le temps le permet ; et
   \item si aucune autre assemblée de mise en candidature n'est tenue, le poste peut être pourvu par une mise en candidature depuis l'assemblée lors de l'assemblée générale annuelle.
  \end{enumerate}
 \item Dès leur mise en candidature, chaque personne candidate doit accomplir les activités préalables de formation de transition pour personnes candidates à l'équipe de direction prescrites par l'équipe de direction avant l'assemblée générale annuelle. Toute personne candidate qui omet de compléter ces activités de formation de transition sera réputée avoir retiré sa candidature et sera remplacée de la manière décrite à la section 4.7 de ces statuts électoraux.
 \item Pour plus de clarté, une personne candidate choisie comme faisant partie de la liste de candidatures ne sera pas considérée comme ayant été élue, nommée ou approuvée au sein de l'équipe de direction tant qu'elle n'aura pas été formellement élue lors de l'assemblée générale annuelle.
 \item Rien dans ces statuts électoraux n'empêchera une mise en candidature faite depuis l'assemblée lors de l'assemblée générale annuelle conformément au paragraphe 163(5) de la Loi canadienne sur les organisations à but non lucratif. Une personne candidate ainsi nominée sera incluse aux côtés de la liste de candidatures pour l'élection lors de l'assemblée générale annuelle.
\end{enumerate}
